\documentclass{article}

\usepackage[margin=1in]{geometry}
\usepackage{graphicx}
\usepackage{xcolor}
\usepackage{titlesec}
\usepackage{enumitem}
\usepackage{hyperref}

% -------------------------------------------------
% Section Styling
\titleformat{\section}
{\color{blue}\normalfont\Large\bfseries}
{\thesection}{1em}{}

% -------------------------------------------------
\title{Group 1 Final Project Weather App}
\author{Group 1}
\date{\today}

\begin{document}

\maketitle

\tableofcontents
\newpage

% -------------------------------------------------
\section{Introduction}

The Group 1 Weather App is a full-stack weather web application that provides users with real-time weather information and a 5-day forecast using the Open-Meteo API. The application was developed using React for the frontend and Node.js with Express for the backend. Docker was used for containerization and Render was used for cloud deployment.

% -------------------------------------------------
\section{Project Overview}

Group 1 Final Project Weather App is a full-stack weather web application built using React, Node.js, and the Open-Meteo API. The application allows users to search for weather conditions by city, view forecasts, manage favorite locations, and track recent searches through a clean and responsive interface.

% -------------------------------------------------
\section{Group Members}

\begin{itemize}[leftmargin=*]
    \item David Gomez
    \item Enrique Cayetano
    \item Saw Rustom Dee
    \item Andres Renteria
    \item Alejandro Villanueva
\end{itemize}

% -------------------------------------------------
\section{Project Goals}

The main goals of this weather application are:

\begin{itemize}[leftmargin=*]
    \item Provide accurate and real-time weather information
    \item Allow users to search weather by city
    \item Display a 5-day weather forecast
    \item Create a responsive and user-friendly interface
    \item Store favorite locations for quick access
    \item Track user search history
    \item Practice full-stack web development concepts
\end{itemize}

% -------------------------------------------------
\section{Purpose of the Project}

Weather applications are commonly used in everyday life to help users plan travel, outdoor activities, and daily schedules. The Group 1 Weather App was created to provide a lightweight and accessible weather platform.

% -------------------------------------------------
\section{Application Preview}

Below is a preview of the weather application interface.

\begin{center}
    \includegraphics[width=0.95\textwidth]{WeatherAppScreenshot.png}
\end{center}

% -------------------------------------------------
\section{Features}

\begin{itemize}[leftmargin=*]
    \item Search weather by city
    \item Display current weather information
    \item Display 5-day weather forecast
    \item Save favorite cities
    \item Remove favorite cities
    \item Store recent search history
    \item Responsive user interface
    \item Real-time weather updates
\end{itemize}

% -------------------------------------------------
\section{API Integration}

This application uses the Open-Meteo API to retrieve:

\begin{itemize}[leftmargin=*]
    \item Current weather conditions
    \item Temperature forecasts
    \item Wind speed data
    \item Humidity information
    \item Daily forecast updates
\end{itemize}

% -------------------------------------------------
\section{System Architecture}

The application follows a full-stack architecture:

\begin{itemize}[leftmargin=*]
    \item Frontend built with React
    \item Backend server built using Node.js and Express
    \item Weather data fetched from Open-Meteo API
    \item Deployment handled through Render
\end{itemize}

% -------------------------------------------------
\section{Technologies Used}

\begin{itemize}[leftmargin=*]
    \item React
    \item Vite
    \item Node.js
    \item Express.js
    \item Open-Meteo API
    \item Docker
    \item Render
    \item GitHub
    \item Jira
    \item TestRail
    \item LaTeX
\end{itemize}

% -------------------------------------------------
\section{Project Structure}

\begin{verbatim}
group-1-final-project/

backend/
    Dockerfile
    package-lock.json
    package.json
    server.js

frontend/
    src/
        App.jsx
        style.css

    Dockerfile
    index.html
    package-lock.json
    package.json

.gitignore
README.md
docker-compose.yml
package.json
\end{verbatim}

% -------------------------------------------------
\section{Installation}

Follow these steps to run the project locally:

\begin{enumerate}[leftmargin=*]

    \item Clone the repository:

\begin{verbatim}
git clone https://github.com/your-username/group-1-final-project.git
\end{verbatim}

    \item Navigate into the project directory:

\begin{verbatim}
cd group-1-final-project
\end{verbatim}

    \item Install dependencies:

\begin{verbatim}
npm install
\end{verbatim}

\end{enumerate}

% -------------------------------------------------
\section{Frontend and Backend Setup}

Open two terminals and run the following commands.

\subsection*{Frontend}

\begin{verbatim}
cd frontend
npm install
npm run dev
\end{verbatim}

\subsection*{Backend}

\begin{verbatim}
cd backend
npm install
npm start
\end{verbatim}

% -------------------------------------------------
\section{Docker Instructions}

The application can also be started using Docker Compose.

\begin{verbatim}
docker compose up --build
\end{verbatim}

After the containers start successfully, open:

\begin{verbatim}
http://localhost:5173
\end{verbatim}

% -------------------------------------------------
\section{Project Link}

Access the deployed project here:

\href{https://group-1-final-project.onrender.com}
{https://group-1-final-project.onrender.com}

% -------------------------------------------------
\section{Usage}

\begin{itemize}[leftmargin=*]
    \item Enter a city name in the search bar
    \item View current weather conditions
    \item Browse the 5-day forecast
    \item Save favorite locations
    \item Review recent search history
\end{itemize}

% -------------------------------------------------
\section{Jira Project Management}

Jira was used to organize tasks, assign responsibilities, and track project progress throughout development.

Jira Link:

\href{https://finalgroup1project.atlassian.net/jira/software/projects/KAN/boards/1?atlOrigin=eyJpIjoiYTMzMTdkZGI1ZTZmNGIxMDgxYWJjZjUzMzRmZmE5YmQiLCJwIjoiaiJ9}
{Jira Project Board}

% -------------------------------------------------
\section{Testing}

The application was tested using TestRail. Test cases and reports were created for Snapshots 2, 3, and 4, along with final system testing.

Testing included:

\begin{itemize}[leftmargin=*]
    \item Frontend functionality
    \item Backend communication
    \item Deployment validation
    \item Docker startup testing
    \item Error handling
\end{itemize}

% -------------------------------------------------
\section{Challenges Faced}

During development, the team encountered several challenges:

\begin{itemize}[leftmargin=*]
    \item Integrating real-time API data
    \item Managing frontend and backend communication
    \item Creating responsive layouts
    \item Handling asynchronous requests
    \item Managing application state
\end{itemize}

% -------------------------------------------------
\section{Learning Outcomes}

Through this project, the team gained experience with:

\begin{itemize}[leftmargin=*]
    \item Full-stack web development
    \item REST API integration
    \item React component architecture
    \item Backend routing with Express
    \item Deployment using Render
    \item Docker containerization
    \item Team collaboration using Git and GitHub
    \item Project management using Jira
    \item Software testing using TestRail
\end{itemize}

% -------------------------------------------------
\section{Future Improvements}

Possible future enhancements include:

\begin{itemize}[leftmargin=*]
    \item Weather maps
    \item Advanced charts and analytics
    \item Push notifications
    \item Additional weather API integrations
    \item Dark mode support
    \item GPS-based location detection
    \item Severe weather alerts
    \item Multi-language support
    \item Temperature unit conversion
    \item User authentication
\end{itemize}

% -------------------------------------------------
\section{License}

This project is for educational purposes.

\end{document}